\section*{Capítulo \ref{ch:b2:crustal-melting} --- La fusión de la corteza: migmatitas y granitos}

\begin{solution}{exo:b2:crustal-melting:1}
(a) Por cada \qty{100}{\gram}: $n(\ce{Al2O3}) = 17.03/101.96 = \qty{0.16703}{mol}$,
$n(\ce{CaO}) = 0.09183$, $n(\ce{Na2O}) = 0.06783$, $n(\ce{K2O}) = 0.03076$.
A/CNK $= 0.16703/0.19042 = 0.877$; A/NK $= 0.16703/0.09859 = 1.694$.
(b) A/CNK $< 1 <$ A/NK: \omterm{def:b2:crustal-melting:asi}{metaluminosa}. El calcio que no pueden admitir los feldespatos va a
silicatos cálcicos: hornblenda (o clinopiroxeno, titanita).
\end{solution}

\begin{solution}{exo:b2:crustal-melting:2}
(a) La moscovita y el granate alojan aluminio más allá de los feldespatos: peraluminoso; este granito
puede ser de tipo S. (b) La hornblenda y la titanita alojan calcio más allá de los feldespatos:
metaluminoso. (c) La egirina y un anfíbol sódico alojan sodio sin el aluminio
del feldespato: peralcalino.
\end{solution}

\begin{solution}{exo:b2:crustal-melting:3}
Un kilómetro de roca pesa $2800\times9.81\times1000 = \qty{2.747e7}{Pa}$, de modo que
$P = 0.2747\,z$ (\unit{kbar}, $z$ en \unit{km}). La recta del solidus es $T = 680 - 7.5(P -
2) = 695 - 2.060\,z$. Igualando $10 + 20z = 695 - 2.060z$ se obtiene $z = 685/22.060 =
\qty{31.1}{km}$, $T = \qty{631}{\celsius}$, $P = \qty{8.5}{kbar}$.
\end{solution}

\begin{solution}{exo:b2:crustal-melting:4}
(a) Las capas blancas son \omterm{def:b2:crustal-melting:migmatite}{leucosomas}; los ribetes negros, \omterm{def:b2:crustal-melting:migmatite}{melanosomas}; el gneis normal, el
mesosoma. (b) \omterm{def:b2:crustal-melting:migmatite}{Leucosomas}, \omterm{def:b2:crustal-melting:migmatite}{melanosomas} y mesosomas discretos y bien definidos: una \omterm{def:b2:crustal-melting:migmatite}{metatexita},
de estructura estromática. (c) La roca se ha convertido en una \omterm{def:b2:crustal-melting:migmatite}{diatexita} (schlierítica y nebulítica):
la fusión ha ido más lejos, los minerales oscuros también se han fundido y el fundido se ha quedado en su sitio,
borrando el bandeado.
\end{solution}

\begin{solution}{exo:b2:crustal-melting:5}
(a) $M = 39.10 + 3\times24.31 + 26.98 + 3\times28.09 + 12\times16.00 + 2\times1.008 =
\qty{417.30}{\gram\per\mole}$; $w = 18.02/417.30 = \qty{4.32}{\percent}$.
(b) $F = 0.25\times0.0432/0.040 = 0.270$, es decir, un \qty{27}{\percent} de fundido en masa: dentro del
20 al \qty{30}{\percent} en volumen producido en experimentos de \omterm{def:b2:crustal-melting:reactions}{fusión por deshidratación} de la biotita con
rocas semejantes.
\end{solution}

\begin{solution}{exo:b2:crustal-melting:6}
(a) Feldespato potásico \qty{278.35}{\gram\per\mole}, albita 262.24, anortita 278.22, moscovita
398.33. (b) Cantidades: feldespato potásico $30/278.35 = \qty{0.10778}{mol}$, albita $0.13347$,
anortita $0.01797$, moscovita $0.01255$. El feldespato potásico es $\tfrac12$\ce{K2O.}$\tfrac12$\ce{Al2O3},
la albita $\tfrac12$\ce{Na2O.}$\tfrac12$\ce{Al2O3}, la anortita \ce{CaO.Al2O3} y la moscovita
$\tfrac12$\ce{K2O.}$\tfrac32$\ce{Al2O3} por unidad de fórmula. $n(\ce{Al2O3}) = 0.05389 +
0.06673 + 0.01797 + 0.01883 = \qty{0.15742}{mol}$; $n(\ce{K2O}) = 0.05389 + 0.00628 =
0.06017$, $n(\ce{Na2O}) = 0.06673$, $n(\ce{CaO}) = 0.01797$, suma $0.14487$.
A/CNK $= 0.15742/0.14487 = 1.087$. (c) Sin la moscovita, los tres feldespatos dan
exactamente 1 (\cref{prop:b2:crustal-melting:feldspar}); la moscovita, con un \ce{Al2O3}
más que el feldespato potásico por unidad de fórmula, hace \omterm{def:b2:crustal-melting:asi}{peraluminosa} la roca.
\end{solution}

\begin{solution}{exo:b2:crustal-melting:7}
Basalto: $900 + 2.5\times(1200 - 850) = \qty{1775}{\mega\joule\per\cubic\metre}$. Corteza:
$2.5\times(850 - 650) + 0.20\times900 = \qty{680}{\mega\joule\per\cubic\metre}$.
$V_c/V_b = 1775/680 = 2.6$. La estimación deja de lado el calor que se pierde por conducción
hacia arriba y hacia los lados y el tiempo que tarda el calor en extenderse; supone además que todo el
calor del basalto llega a una corteza lista para fundirse: es una cota superior.
\end{solution}

\begin{solution}{exo:b2:crustal-melting:8}
(a) $n(\ce{Al2O3}) = 14.15/101.96 = 0.13878$, $n(\ce{CaO}) = 0.01753$, $n(\ce{Na2O}) =
0.07809$, $n(\ce{K2O}) = 0.04416$; A/CNK $= 0.13878/0.13978 = 0.993$, A/NK $=
0.13878/0.12225 = 1.135$: \omterm{def:b2:crustal-melting:asi}{metaluminosa}, por poco. (b) Con un \qty{14.30}{\percent} de \ce{Al2O3},
$0.14025/0.13978 = 1.003$; con un \qty{14.00}{\percent}, $0.13731/0.13978 = 0.982$. El índice
queda a ambos lados de 1 dentro de la incertidumbre certificada: la obsidiana está en el
límite entre \omterm{def:b2:crustal-melting:asi}{metaluminosa} y \omterm{def:b2:crustal-melting:asi}{peraluminosa}, y su clase no es robusta.
\end{solution}

\begin{solution}{exo:b2:crustal-melting:9}
(a) Medias $\bar P = \qty{5.0}{kbar}$, $\bar w = \qty{10.045}{\percent}$;
$\sum(P - \bar P)(w - \bar w) = 17.62$, $\sum(P - \bar P)^2 = 14.0$; pendiente
$17.62/14.0 = \qty{1.259}{\percent}$ por kilobar, ordenada en el origen $10.045 - 1.259\times5.0 =
\qty{3.75}{\percent}$. La pendiente coincide con el 1.27 citado con una diferencia inferior al
\qty{1}{\percent}. (b) $3.75 + 1.259\times6 = \qty{11.30}{\percent}$. (c) $10/0.1130 =
\qty{88.5}{\gram}$ de fundido. (d) La recta se apoya en datos entre 3 y \qty{8}{kbar} a
\qty{800}{\celsius} solamente; nada garantiza que la solubilidad siga siendo lineal a presiones más
altas, donde el fundido y el fluido se parecen cada vez más.
\end{solution}

\begin{solution}{exo:b2:crustal-melting:10}
(a) La recta sube \qty{4}{kbar} en \qty{50}{\celsius}: $P = 6 + 0.08\,(T - 750)$;
a \qty{770}{\celsius}, $P = \qty{7.6}{kbar}$. (b) De 7.6 a \qty{6}{kbar} la
temperatura baja a lo largo de la recta de 770 a \qty{750}{\celsius}; el calor liberado,
$2.5\times20 = \qty{50}{\mega\joule\per\cubic\metre}$, funde $50/900 = 0.056$, cerca de un
\qty{6}{\percent} de la roca en volumen. (c) Es la ventana propuesta para los
\omterm{def:b2:crustal-melting:leucogranite}{leucogranitos} del Himalaya, de 6 a \qty{8}{kbar} y de 750 a \qty{770}{\celsius}: la descompresión
por sí sola produce fundido, pero solo un pequeño tanto por ciento en cada cruce de la curva.
\end{solution}

\begin{solution}{exo:b2:crustal-melting:11}
$\ell^2/\kappa = (2000)^2/10^{-6} = \qty{4e12}{s} = \qty{0.13}{Myr}$ para el sill;
$(10^4)^2/10^{-6} = \qty{1e14}{s} = \qty{3.2}{Myr}$ para la capa subplacada. Ambos son cortos
frente a diez millones de años: un solo episodio no puede mantener fundida la corteza tanto tiempo, lo que requiere
un aporte de calor renovado una y otra vez, por intrusiones repetidas o por el lento calentamiento de una
corteza engrosada.
\end{solution}

\begin{solution}{exo:b2:crustal-melting:12}
(a) $5.0x + 2.0(1 - x) = 2.75$ da $x = 0.75/3.0 = 0.25$: una cuarta parte del neosoma es
\omterm{def:b2:crustal-melting:migmatite}{leucosoma}. (b) $0.25\times0.20 + 0.75\times4.0 = 0.05 + 3.00 = \qty{3.05}{\percent}$, como
en el mesosoma: la interpretación en sistema cerrado se sostiene. (c) Por cada \qty{100}{\gram} de
\omterm{def:b2:metamorphic-facies:metamorphism}{protolito} se formaron \qty{25}{\gram} de fundido y \qty{75}{\gram} de residuo; si se fueron $m$ gramos de
fundido, $(25 - m)/(100 - m) = 0.10$, de modo que $25 - m = 10 - 0.1m$ y $m = 15/0.9 =
\qty{16.7}{\gram}$: dos tercios del fundido abandonaron el afloramiento.
\end{solution}

\begin{solution}{pb:b2:crustal-melting:1}
\noindent\textbf{Parte I.}

\textbf{1.} $39.10 + 3\times26.98 + 3\times28.09 + 12\times16.00 + 2\times1.008 =
\qty{398.33}{\gram\per\mole}$.

\textbf{2.} $18.02/398.33 = \qty{4.52}{\percent}$.

\textbf{3.} $200\times0.0452 = \qty{9.04}{\gram}$ de agua por kilogramo de esquisto.

\textbf{4.} $9.97 + (13.82 - 9.97)\times2/3 = \qty{12.54}{\percent}$. El fundido mínimo es
$9.04/0.1254 = \qty{72.1}{\gram}$ por kilogramo: $F \ge \qty{7.21}{\percent}$
(\cref{prop:b2:crustal-melting:water}).

\textbf{5.} $9.04/0.060 = \qty{151}{\gram}$: $F = \qty{15.1}{\percent}$.

\textbf{6.} Por ejemplo: la roca puede contener menos moscovita; parte de la moscovita puede
haberse descompuesto ya por debajo del solidus, con el cuarzo, en feldespato potásico y
aluminosilicato, y su agua haberse ido como fluido antes de que empezara la fusión; el fundido puede contener más
agua que un \qty{6}{\percent}; algo de agua puede quedarse en la biotita.

\noindent\textbf{Parte II.}

\textbf{7.} $680 - 7.5\times(7.8 - 2) = \qty{636.5}{\celsius}$; la roca alcanza su pico
\qty{43.5}{\celsius} por debajo y no puede fundir, ni siquiera con agua libre.

\textbf{8.} $605 + (19.2 - 15.3)/(27 - 15.3)\times65 = \qty{626.7}{\celsius}$. En
$t = 0$ la base, a \qty{595}{\celsius}, está \qty{32}{\celsius} por debajo; a los
\qty{10}{Myr}, a \qty{656}{\celsius}, está por encima.

\textbf{9.} La fusión en el solidus húmedo necesita agua libre, y una roca metamórfica calentada
contiene muy poca: el primer fundido consume el agua libre y deja de crecer
hasta que se descomponen las micas.

\textbf{10.} Entre 30 y \qty{40}{Myr}: $30 + (800 - 795)/(845 - 795)\times10 =
\qty{31.0}{Myr}$.

\textbf{11.} La base baja de \qty{12.5}{kbar} entre 70 y \qty{80}{Myr}:
$70 + (13.7 - 12.5)/(13.7 - 12.4)\times10 = \qty{79.2}{Myr}$, a
$914 - 14\times0.923 = \qty{901}{\celsius}$, dentro del intervalo del anfíbol: una anfibolita
en la base estaría fundiéndose.

\textbf{12.} $90 + (11.0 - 10)/(11.0 - 9.6)\times10 = \qty{97.1}{Myr}$, a $867 -
53\times0.714 = \qty{829}{\celsius}$. Eso está por encima del intervalo de la moscovita: la moscovita
se ha descompuesto. Está dentro de los 800 a \qty{850}{\celsius} de la descomposición de la biotita
(medidos a menor presión): la biotita, el otro mineral hidratado del esquisto, está en su
límite.

\noindent\textbf{Parte III.}

\textbf{13.} $n(\ce{Al2O3}) = 15.0/101.96 = \qty{0.14712}{mol}$, $n(\ce{CaO}) =
0.55/56.08 = 0.00981$, $n(\ce{Na2O}) = 3.4/61.98 = 0.05486$, $n(\ce{K2O}) = 4.9/94.20 =
0.05202$.

\textbf{14.} A/CNK $= 0.14712/0.11669 = 1.261$.

\textbf{15.} A/NK $= 0.14712/0.10688 = 1.377$; A/CNK $> 1$: peraluminoso.

\textbf{16.} Peraluminoso y procedente de un esquisto: un \omterm{def:b2:crustal-melting:types}{granito de tipo S}. Su exceso de aluminio se aloja en
la moscovita y el granate (o la cordierita).

\textbf{17.} $0.14712 - 0.11669 = \qty{0.03043}{mol}$ de \ce{Al2O3}; en forma de moscovita,
$0.03043\times398.33 = \qty{12.1}{\gram}$ en \qty{100}{\gram} de granito.

\textbf{18.} La obsidiana es (por poco) \omterm{def:b2:crustal-melting:asi}{metaluminosa}, sin el exceso de aluminio de un
fundido metasedimentario: su fuente es probablemente distinta. Pero una sola muestra prueba poco:
toda serie de granitoides tiene miembros ligeramente peraluminosos, y una serie se juzga por toda su
tendencia.

\noindent\textbf{Parte IV.}

\textbf{19.} $10^{12}\,\unit{\cubic\metre}\times2650 = \qty{2.65e15}{kg}$.

\textbf{20.} El granito es el fundido menos su agua: $2.65\times10^{15}/0.94 =
\qty{2.82e15}{kg}$ de fundido.

\textbf{21.} Solo $0.151\times0.70$ de la masa del esquisto sale como fundido:
$2.82\times10^{15}/(0.151\times0.70) = \qty{2.67e16}{kg}$.

\textbf{22.} $2.67\times10^{16}/2800 = \qty{9.5e12}{\cubic\metre}$, es decir,
\qty{9.5e3}{\cubic\kilo\metre}, casi diez veces el \omterm{def:b2:magma-emplacement:bodies}{plutón}.

\textbf{23.} $9.5\times10^3/5 = \qty{1.9e3}{\square\kilo\metre}$, un cuadrado de unos
\qty{44}{km} de lado.

\textbf{24.} Con $F = \qty{7.21}{\percent}$, la fuente es $15.1/7.21 = 2.1$ veces
mayor, unos \qty{2.0e4}{\cubic\kilo\metre}. El volumen de corteza fuente fundida para formar
el \omterm{def:b2:magma-emplacement:bodies}{plutón}, con la \omterm{def:b2:mantle-melting:fraction}{fracción de fundido} de la pregunta~5, es
\textbf{\qty[text-series-to-math=true]{9.5e3}{\cubic\kilo\metre}}.
\end{solution}
